\section{Baseline Reproduction as a Validity Check}
\label{sec:baselinecheck}

The companion article's equivalence invariant~\cite{methodpaper} asserts
that the framework with all modules disabled \emph{is} the published
pipeline. That is a claim about code
structure, so measurement cannot establish it --- but measurement can
falsify it, and we checked two ways.

The Baseline arm's success rate reproduces the original combined-stress
figure across three measurements: 52.1\% in the original failure-mode
campaign, 52.1\% in the ablation reported here, and 51.1\% in an interim
run at 139 of 144 trials.\footnote{The interim analysis predates the five
heaviest configurations finishing. The full 144-trial results reported
throughout supersede it; conclusions are unchanged.} The first is not a
re-run of the same function: the original campaign drove the published
pipeline through its own entry point, whereas the Baseline arm is the
adaptive entry point with all flags disabled. Their agreement is
therefore evidence \emph{across} implementations, which a single-function
re-run could not provide.

The second check is stronger than agreement between rates. The
instrumented re-run of Sec.~\ref{sec:design} was gated on reproducing the
recorded per-trial outcomes of \emph{every} arm exactly --- trial by
trial, not in aggregate --- and did so on all 144 trials in all five arms.
Adding the instrumentation that Sec.~\ref{sec:matched} depends on
therefore changed no result, and the seed-paired determinism the whole
design rests on holds in practice rather than merely in principle. Had it
not, the matched-resource controls would have been uninterpretable, since
each is built from a diagnostic recorded during that run.
